\chapter{Introduction}
\label{chap:introduction}

\chapterintro{%
Trusted-relay quantum key distribution (QKD) networks must coordinate quantum-link
resources with conventional network control. This thesis addresses one recurring
control-plane computation: evaluating a bounded set of admissible route candidates
from key-resource and network-state metadata. It develops \QFlow, a deterministic
fixed-point FPGA accelerator, and evaluates five independently implemented hardware
configurations under a common physical measurement contract.}

\section{Motivation}

\subsection{QKD links and trusted-relay networking}

QKD enables two endpoints to establish shared secret material by combining quantum
transmission with authenticated classical communication. After sifting, parameter
estimation, error correction, and privacy amplification, the usable keys are
classical data managed by a key management system (KMS). A network service therefore
depends not only on the quantum link but also on key storage, authorization,
consumption, replenishment, monitoring, and failure handling
\cite{cao2022qinternet,mehic2020networking}.

Direct QKD links do not necessarily connect every communicating pair. Trusted-relay
networks extend coverage by forwarding key material or coordinating secure segments
through intermediate trusted sites. Operational demonstrations such as MadQCI and a
carrier-scale network show the importance of integrating QKD devices with KMS and
software-defined networking (SDN) functions
\cite{martin2024madqci,chen2025carrier}. Route selection becomes a control problem
because candidate paths can differ in distance, hop count, available key material,
key-generation rate, utilization, and operational health.

\begin{figure}[H]
\centering
\resizebox{\textwidth}{!}{%
\begin{tikzpicture}[
  ext/.style={draw=QFlowGray,fill=gray!7,rounded corners=2pt,
    minimum height=1.15cm,text width=2.35cm,align=center,font=\small},
  core/.style={draw=QFlowBlue,fill=QFlowLightBlue,rounded corners=2pt,
    minimum height=1.15cm,text width=2.45cm,align=center,font=\small\bfseries},
  arr/.style={-{Latex[length=2.3mm]},thick,draw=QFlowNavy},
  node distance=0.55cm]
\node[ext] (qkd) {QKD devices and link monitors};
\node[ext,right=of qkd] (kms) {KMS and SDN controller};
\node[ext,right=of kms] (gen) {candidate generation};
\node[core,right=of gen] (eval) {\QFlow\\candidate evaluation};
\node[ext,right=of eval] (install) {policy validation and route installation};
\draw[arr] (qkd)--(kms);
\draw[arr] (kms)--(gen);
\draw[arr] (gen)--(eval);
\draw[arr] (eval)--(install);
\end{tikzpicture}}
\caption{System context of \QFlow. Candidate generation, hard-eligibility
computation, and route installation are control-plane functions outside the FPGA
kernel.}
\label{fig:system-context}
\end{figure}

\subsection{Key-pool-aware routing control}

Post-processed keys are finite classical resources. For a directed link, the KMS can
report pool occupancy, effective key-generation rate, consumption, link status,
utilization, and externally measured health indicators such as QBER-derived status.
An optional age or expiry field represents KMS or security policy; it is not a
quantum coherence property of stored bits. A route that is short in distance may be
undesirable when one of its key pools is depleted, while a slightly longer route may
provide greater reserve or replenishment capacity. Recent work consequently studies
dynamic key-aware, multipath, joint-resource, and reinforcement-learning approaches
\cite{bi2023dynamic,kiktenko2024nonoverlap,zhang2024rckta,hao2025qlearning,
johann2026adaptive}.

Two decision layers are useful. Hard eligibility first excludes a route whose links
are unavailable, under-provisioned, or administratively invalid. Soft preference
then ranks the remaining routes according to a selected scoring profile. Keeping
these layers separate prevents a high score from making an inadmissible route
appear usable.

\subsection{Route-candidate evaluation as the selected problem}

Full graph routing includes topology acquisition, constraint processing, candidate
generation, global optimization, security validation, and route installation.
Mapping that entire workflow to a small FPGA would mix irregular control tasks with
a compact recurring arithmetic task. This thesis instead accelerates the latter.
For source $s$ and destination $d$, software supplies
\begin{equation}
\mathcal{R}_{sd}
=
\{\mathcal{R}_0,\mathcal{R}_1,\ldots,\mathcal{R}_{m-1}\}.
\label{eq:intro-candidate-set}
\end{equation}
Each candidate $\mathcal{R}_i$ has a fixed-width descriptor containing an upstream
validity result, path cost, a bounded route-quality code, utilization, hop count,
and a deterministic identifier. \QFlow\ selects a scoring profile, gates candidates
using the supplied validity bits, evaluates the remaining descriptors, and returns
the chosen identifier, score, profile, and status.

This boundary makes exact verification possible. The same descriptor and fixed-point
contract can be evaluated in software, RTL simulation, and physical hardware. It
also reflects a practical heterogeneous design in which the CPU retains graph and
policy responsibilities while the FPGA offloads repeated deterministic evaluation.

\subsection{Why deterministic hardware acceleration is relevant}

A candidate evaluator applies the same bounded operations to every request:
quantization, read-only policy lookup, supplied-validity gating, weighted fixed-point
scoring, comparison, and deterministic tie-breaking. These operations are suitable
for dedicated pipelines and finite-state control. FPGA implementation provides
cycle-level predictability, explicit numeric widths, and reconfiguration as the
controller or scoring profiles evolve.

The present research stage also benefits from avoiding ASIC fabrication and from
implementing all comparison configurations on the same device. CPU software remains
essential for orchestration, candidate generation, offline learning, and irregular
tasks. The platform choice is therefore a partitioning decision rather than a claim
that one platform is universally superior.

\section{Research Gap}

The literature contains sophisticated software studies of QKD routing and resource
allocation, including adaptive and reinforcement-learning methods. Separately,
FPGA-based QKD research has demonstrated high-throughput reconciliation,
post-processing, and synchronization
\cite{lu2023multidimensional,venkatachalam2023postprocessing,liao2025polar,
chandra2025synchronization}. Within the documented literature review, no directly
comparable same-board implementation of the evaluated candidate-selection function
was identified. The gap is therefore an implementation and evidence gap: a compact,
adaptive, fixed-point candidate evaluator has not been compared through independent
physical configurations with a common interface, exact reference agreement, and
post-route timing and resource evidence.

\section{System Scope and Boundary}

\QFlow\ is a route-candidate evaluator. It consumes candidate descriptors and
network-state metadata that have already been validated and quantized at the
software/hardware boundary. It supports fixed, rule-adaptive, and offline-learned
profile selection. The learned configuration performs online inference from a
read-only policy; Q-values are not updated during FPGA execution.

Global topology management, candidate generation, KMS operation, SDN orchestration,
security-policy enforcement, secret-key delivery, and route installation remain in
software or the control plane. The FPGA is neither a complete graph router nor a QKD
physical-layer device. The physical experiment measures the accepted-request-to-done
kernel boundary; it does not measure complete network provisioning latency.

\section{Problem Statement}

Given an externally generated candidate set $\mathcal{R}_{sd}$ and quantized
key-resource and network-state metadata, the research problem is to design a bounded
hardware kernel that:

\begin{enumerate}
  \item excludes candidates whose upstream validity bits report failed hard
  constraints;
  \item selects a fixed or adaptive scoring profile deterministically;
  \item evaluates feasible candidates using reproducible fixed-point arithmetic;
  \item returns a stable result under a synchronous request/done protocol; and
  \item is physically implementable on a resource-constrained FPGA with measurable
  timing, resource, exactness, and adaptive-behavior trade-offs.
\end{enumerate}

\section{Research Objectives}

The objectives are to:

\begin{enumerate}
  \item formulate a classical key-resource model based on pool occupancy,
  generation rate, consumption, status, utilization, health indicators, and
  policy-controlled freshness;
  \item define a fixed-point candidate-scoring contract with distinct upstream hard
  eligibility, FPGA validity gating, and soft preference;
  \item develop fixed, rule-adaptive, and offline-learned profile controllers;
  \item implement five independent FPGA configurations under one external contract;
  \item verify reference-to-RTL and reference-to-board exactness;
  \item measure post-route timing, kernel latency, resources, and profile behavior;
  and
  \item interpret route-quality and platform trade-offs without crossing evidence
  boundaries.
\end{enumerate}

\section{Research Questions}

\begin{description}[leftmargin=1.3cm,style=nextline]
  \item[\textbf{RQ1.}] Can all five independently implemented configurations meet a
  100~MHz post-route target on the selected FPGA?
  \item[\textbf{RQ2.}] Do physical outputs match the fixed-point reference exactly
  under a common replay and output schema?
  \item[\textbf{RQ3.}] What latency, maximum-frequency, and resource trade-offs
  arise across the fixed and adaptive configurations?
  \item[\textbf{RQ4.}] Does the offline-learned controller exhibit distinct and
  controlled profile behavior relative to the rule-adaptive controller?
  \item[\textbf{RQ5.}] Do the held-out physical comparisons establish statistically
  significant route-quality improvement for the learned configuration?
\end{description}

\section{Research Hypotheses}

\begin{description}[leftmargin=1.3cm,style=nextline]
  \item[\textbf{RH1.}] Each independent hardware evaluation configuration will meet
  the 100~MHz post-route timing target.
  \item[\textbf{RH2.}] Every checked physical result field will match the
  registered fixed-point reference.
  \item[\textbf{RH3.}] The reinforcement-learned adaptive configuration will retain
  microsecond-scale kernel latency with controlled implementation overhead.
  \item[\textbf{RH4.}] The reinforcement-learned adaptive configuration with
  minimum-dwell control will produce fewer profile transitions than the
  rule-adaptive configuration on the common replay.
  \item[\textbf{RH5.}] The reinforcement-learned adaptive configuration will achieve
  statistically significant route-quality improvement over fixed-profile and
  rule-adaptive alternatives on the held-out physical pairs.
\end{description}

\section{Contributions}

The thesis makes the following contributions:

\begin{enumerate}
  \item a classical key-resource model for candidate evaluation based on pool
  occupancy, generation rate, consumption, operational status, policy freshness,
  and externally measured health indicators;
  \item a deterministic fixed-point candidate-evaluation contract that separates
  upstream hard-eligibility computation from FPGA validity gating and
  profile-conditioned preference;
  \item an adaptive controller with four radix-4 features, 256 states, four full
  software-training profile actions, offline tabular Q-learning, and a 256-entry
  two-bit policy ROM whose physical candidate shell deploys the profiles' objective
  ratios with minimum-dwell control;
  \item five independent FPGA implementations with a common interface and
  measurement boundary;
  \item a same-board campaign comprising 420 physical records, 3,780 checked fields,
  and exact reference agreement; and
  \item a complete engineering comparison of timing, latency, resources, adaptive
  behavior, route-quality statistics, and stage-specific platform implications.
\end{enumerate}

\section{Practical Significance}

In a prospective KMS/SDN deployment, candidate evaluation may recur whenever
requests arrive or resource state changes. Offloading this fixed-width calculation
can provide predictable control latency and a bit-exact implementation shared across
tests and deployments. Reprogrammable policy and profile storage also allows the
accelerator to evolve without replacing the surrounding control plane.

\section{User Requirements and Broader Considerations}

The engineering users of the proposed kernel are network-control and security
software developers. Their principal requirements are deterministic requests and
responses, explicit numeric formats, visible failure status, reproducible profiles,
and a software fallback. A deployment must authenticate the host/FPGA interface,
validate inputs, protect policy artifacts, log configuration identity, and retain
the controller's authority over final route installation.

The work does not process human-subject data. Its societal relevance lies in
supporting reproducible research on secure-network control rather than promising a
complete commercial service. Hardware safety is limited to normal laboratory
electrical and electrostatic-discharge practice. Environmental benefit is not
quantified because no activity-calibrated power measurement or life-cycle study was
performed; reconfigurability can, however, extend prototype usefulness by allowing
new bitstreams on the same device.

\section{Thesis Organization}

Chapter~\ref{chap:background} reviews trusted-relay QKD networking, key-resource
control, multi-objective scoring, reinforcement learning, fixed-point arithmetic,
platform alternatives, and related work. Chapter~\ref{chap:methodology} defines the
system model, scoring profiles, hardware configurations, controller, RTL,
verification, FPGA experiment, statistical method, and supporting VLSI method.
Chapter~\ref{chap:results} presents and discusses the software, RTL, post-route,
physical replay, adaptive-behavior, route-quality, platform, and VLSI results.
Chapter~\ref{chap:conclusion} summarizes the findings, contributions, limitations,
and future work. The appendices provide reproducibility records, fixed-point and
protocol details, additional results, and experimental-scope notes.

\chapterconclusion{%
The selected problem is a bounded, recurring route-candidate evaluation task within
a larger KMS/SDN workflow. The research therefore concentrates on fixed-point
correctness, deterministic control, physical FPGA implementation, and transparent
trade-offs. The next chapter establishes the technical and literature foundations
for that design.}
